% Generated draft: requires mathematical and source review.
\documentclass[11pt]{article}
\usepackage[margin=25mm]{geometry}
\usepackage{fontspec}
\setmainfont{Latin Modern Roman}
\usepackage{amsmath,amssymb,mathtools,unicode-math}
\setmathfont{Latin Modern Math}
\usepackage{enumitem,needspace,xurl,hyperref}
\setlength{\parindent}{0pt}
\setlength{\parskip}{5pt}
\setlength{\emergencystretch}{3em}
\newcommand{\N}{\mathbb{N}}
\newcommand{\Z}{\mathbb{Z}}
\newcommand{\Q}{\mathbb{Q}}
\newcommand{\R}{\mathbb{R}}
\newcommand{\C}{\mathbb{C}}
\newcommand{\F}{\mathbb{F}}
\newcommand{\A}{\mathbb{A}}
\newcommand{\PP}{\mathbb{P}}
\newcommand{\E}{\mathbb{E}}
\newcommand{\eps}{\varepsilon}
\newcommand{\dd}{\,\mathrm{d}}
\begin{document}
\section*{UnsolvedMath ID: 30005666}
\textbf{Source problem code:} catalogue entry
\subsection{Definitions and mathematical statement}
The supplied registry asks: \emph{Which subgraphs of the curve complex are rigid?} The surface, the meaning of rigidity, and even the antecedent of ``the complex'' in the extracted passage are unresolved.

For a conventional interpretation, let \(S\) be a connected orientable surface of finite type. An essential, nonperipheral simple closed curve is an embedded circle that bounds neither a disk nor a once-punctured disk and is not isotopic to a boundary component. The usual disjointness curve graph \(\mathcal C^1(S)\) has isotopy classes of such curves as vertices; distinct vertices are adjacent when they admit disjoint representatives. The corresponding curve complex has simplices given by collections admitting pairwise disjoint representatives. These are conditional conventions: the record specifies neither this surface class nor whether modified adjacency conventions for surfaces of low complexity are intended.

One possible precise framework for the requested classification is the following. Fix \(S\), a collection \(\mathscr H\) of subgraphs of \(\mathcal C^1(S)\), a group \(G\) acting on \(\mathcal C^1(S)\), and, for each \(H\in\mathscr H\), a specified class \(\mathscr M(H,\mathcal C^1(S))\) of admissible maps. Define the extension property by
\[
\operatorname{Ext}_{G,\mathscr M}(H)
\quad\Longleftrightarrow\quad
\forall f\in\mathscr M(H,\mathcal C^1(S))\;
\exists g\in G\;
\forall v\in V(H),\quad f(v)=g\cdot v.
\]
Under an extension-based interpretation of rigidity, the problem is to characterize
\[
\bigl\{H\in\mathscr H:\operatorname{Ext}_{G,\mathscr M}(H)\bigr\}.
\]
The supplied record does not determine \(\mathscr H\), \(G\), or \(\mathscr M\), or whether uniqueness of an extension is required. In particular, it does not specify finite versus arbitrary subgraphs, induced versus arbitrary subgraphs, locally injective simplicial maps versus other maps, or graph automorphisms versus homeomorphism-induced actions. Consequently, this is a conditional formalization schema, not an identification of the intended theorem. The source context must be recovered before a unique mathematical statement can be fixed.
\subsection{Short English statement}
Characterize the rigid subgraphs of the curve complex. The supplied extract attributes a related question to McCarthy--Papadopoulos, but leaves the surface, the meaning of rigidity, and the referent of ``the complex'' unclear. The displayed extension property describes a possible formulation whose parameters require source confirmation.
\subsection{Sources}
\begin{itemize}
\item Local UnsolvedMath registry: \texttt{data/unsolvedmath/problems.json}, permanent ID \(30005666\), title \emph{Rigid Subgraphs of the Curve Complex}, source problem code \texttt{catalogue entry}. The supplied record labels the statement \texttt{corrected\_verified} and reports direct checking of the official report; this is supplied provenance, not independent verification by this draft.
\item Supplied report citation: \emph{Teichmüller Theory: Classical, Higher, Super and Quantum} (2024), Oberwolfach Reports, pp.~1843--1892, DOI \texttt{10.4171/owr/2023/33}. \url{https://doi.org/10.4171/owr/2023/33}. The record separately gives the proposed year as 2023. The precise page containing the problem is not supplied.
\item The extracted passage attributes the question concerning subgraphs of an antecedent complex to McCarthy--Papadopoulos and describes that complex as a natural counterexample to Ivanov's metaconjecture. It then runs into a contribution headed \emph{The curve complex via model theory}, credited to Valentina Disarlo, jointly with Javier de la Nuez Gonzalez and Thomas Koberda. This adjoining heading does not establish the referent of the preceding passage or authorship of the problem.
\item Aramayona--Leininger, \emph{Finite rigid sets in curve complexes}. \url{https://arxiv.org/abs/1206.3114}.
\item \emph{Finite rigid sets and the separating curve complex}. \url{https://doi.org/10.1016/j.topol.2022.108078}.
\item Supplied literature assessment, checked \textbf{2026-08-22}: status \texttt{partially\_solved}. The supplied assessment states that finite rigid sets are known for the full curve complex and several subcomplexes, while a general classification of the subgraphs requested in the extracted passage is unavailable. This is supplied dated curation, not an independently verified literature assessment; the background explicitly requests independent verification of time-sensitive status and cited claims.
\end{itemize}
\end{document}
